(1) 均匀球在外静电场中的内外 Laplace 边界匹配给感应偶极矩
\[
p=\alpha_eE_0,\quad
\alpha_e=4\pi\epsilon_0a^3\frac{\epsilon_r-1}{\epsilon_r+2}.
\]
该式要求球相对于波长足够小，使其整个体积近似处于同一相位的电场。

(2) 频率 \(\omega=ck\) 的谐振电偶极平均辐射功率
\[
P_{\rm rad}=\frac{\omega^4|p|^2}{12\pi\epsilon_0c^3}.
\]
入射平面波平均强度 \(I_{\rm in}=\epsilon_0c|E_0|^2/2\)。散射截面定义为 \(\sigma_s=P_{\rm rad}/I_{\rm in}\)，代入并约去 E0 得
\[
\sigma_s=\frac{k^4|\alpha_e|^2}{6\pi\epsilon_0^2}
=\frac{8\pi}{3}k^4a^6
\left|\frac{\epsilon_r-1}{\epsilon_r+2}\right|^2.
\]

(3) 对固定材料，\(\sigma_s\propto a^6/\lambda^4\)，体现小球尺寸的高次敏感性以及短波长更强的 Rayleigh 散射。此散射截面可以小于也可以大于几何截面，但在 ka≪1 时比值通常按 \((ka)^4\) 很小。

(4) 当 \(\epsilon_r\to1\)，球与真空不可区分，\(\alpha_e\to0,\sigma_s\to0\)。当 ka 不再远小于1，球内相位不均匀、磁偶极及更高多极激发不能忽略，需要采用完整 Mie 散射理论。
